\textbf{Data and model.} The frozen model was FasterDecoding/medusa-vicuna-7b-v1.3 on one 16-GiB RTX 3080 Ti Laptop GPU. The loader quantized the backbone to NF4 from FP16 and promoted floating-point execution and the five prediction heads to FP32, with TF32 disabled. PyTorch reported 6.93 GiB of allocated GPU tensors after loading. Peak allocation during the seven instrumented diagnostic prompts ranged from 9.51 to 9.58 GiB. All policies shared the same weights, top ten candidates per head and 64-token generation cap. We used the original AlpacaEval split and MT-Bench first turns. Training prompt 801 duplicated MT-Bench question 81 verbatim. Removing the training copy before fitting left 199 training, 50 validation and 80 test inputs. A Unicode-, case- and whitespace-normalized text check found no duplicates among the remaining inputs.

\textbf{Reference trajectory.} At each ordinary cached forward, we recorded the backbone's greedy choice, its probability, and the ranked token IDs and probabilities from the five heads. The next backbone token became the input to the next forward. For head $j$ at state $t$, the label is the rank of the $j$th token after the backbone choice, or a miss if that token falls outside the top ten. Thus the label and candidates refer to one recorded trajectory. This avoids deriving features from a separate teacher-forced execution. An independent audit reconstructs every stored label directly from the ranked candidates and generated IDs. Live tree verification is checked separately.

Training yielded 10,854 complete five-head horizons on 196 prompts. Validation yielded 2,748 on 49 prompts. Three training answers and one validation answer reached EOS before a full horizon and supplied zero labeled states. The coverage manifest retains all four. Only positions with all five future targets enter replay. The endpoint therefore describes interior states of bounded generations.

\textbf{Model selection.} The integrated calibration notebook tunes count smoothing and logistic regularization and rank cap on the validation prompts. It compares three candidate-pool constructions for the direct classifier by validation acceptance at $B=63$. The common calibration pool uses MEDUSA+ to select 64 candidate paths per sampled state (up to 4,000 states per split). Brier score averages squared path-probability errors. Expected calibration error (ECE) uses 15 fixed-width probability bins. A grouped feature ablation refits the conditional classifier with each feature group omitted. The frozen notebook fixes the seed and comparison rules before the evaluation results are used.

\textbf{Endpoints.} Replay scores all five trees at each eligible evaluation state and weights states equally. We use 2,000 paired prompt-bootstrap resamples for differences, retaining all states from each sampled prompt. The fixed contrast uses the marginal tree as its baseline. The dynamic contrast uses MEDUSA+. These are per-budget intervals. The analysis treats the family of budget comparisons as exploratory.

\textbf{Live verification and timing.} We fixed an eight-question panel before fitting by selecting the first MT-Bench question from each category, with a separate writing warmup. For budgets 16 and 63 we check cached-reference token IDs, final output IDs, round prefixes, full candidate ranks and selected-node accept/reject decisions. A separate timing pass runs the five policies and cached greedy in seeded, randomized order within prompts. Throughput divides total generated tokens by total measured time, including prefill and generation. Model loading, tokenization, cache allocation and fixed-tree buffer setup lie outside the timer. A separate synchronized run profiles tree overhead on the first measured prompt.
